\chapter{Introduction}

This blueprint documents a Lean~4 / Mathlib formalization of
\emph{Integer Coefficients Power Series with Prescribed Zero Sets} by
Jon Bannon and David Feldman. The paper proves that a discrete effective
divisor on the open unit disk $\D$ is the zero divisor of a holomorphic
function on $\D$ with integer Taylor coefficients if and only if it is
invariant under complex conjugation (Theorem~\ref{thm:main}), via an
intermediate result realizing every effective divisor with Gaussian-integer
coefficients (Theorem~\ref{thm:propZi}).

The construction in both theorems deforms the classical Weierstrass
elementary factor by a free parameter, so that each modified factor of
order $n$ leaves all Taylor coefficients of degree $\le n$ unchanged while
shifting the coefficient of degree $n+1$ by a controlled affine amount.
This \emph{triangular} structure lets the coefficients of an infinite
product be rounded to integers one at a time, without disturbing
previously-fixed coefficients or breaking convergence.

\bigskip
\noindent\textbf{Scope.} This blueprint covers the paper's
Sections~1--3 (the modified elementary factors, the affine-control lemma,
and the proofs of Theorems~\ref{thm:propZi} and~\ref{thm:main}). The
paper's remaining sections (the worked numerical example of
Section~4, and the ring-theoretic consequences of Section~5 — the
characterization of units, the ideal-theoretic corollaries, and the
comparison of $\mathrm{Spec}$ and $\mathrm{MaxSpec}$) are \emph{not}
formalized here.

\bigskip
\noindent\textbf{How to read this blueprint.} Every definition and
theorem below is annotated with the Lean declaration(s) that formalize it
(via a \lean{} link) and a green checkmark once fully proved
(\leanok). The dependency graph (linked from the web version's
navigation) shows how each result builds on the others. All results in
this blueprint are fully proved in Lean, with no \texttt{sorry} and no
axioms beyond Mathlib's standard three (\texttt{propext},
\texttt{Classical.choice}, \texttt{Quot.sound}).

\chapter{Preliminaries}

\begin{definition}[The open unit disk]
  \label{def:disk}
  \lean{Weierstrass.𝔻}
  \leanok
  $\D = \{z \in \C : \abs{z} < 1\}$.
\end{definition}

\begin{definition}[Holomorphic on $\D$]
  \label{def:holomorphicOn}
  \uses{def:disk}
  \lean{Weierstrass.HolomorphicOn}
  \leanok
  A function $f : \C \to \C$ is \emph{holomorphic on $\D$} if it is
  complex-analytic at every point of $\D$.
\end{definition}

\begin{definition}[Taylor coefficient]
  \label{def:taylorCoeff}
  \lean{Weierstrass.taylorCoeff}
  \leanok
  The $n$-th Taylor coefficient of $f$ at the origin is
  $[z^n]f := f^{(n)}(0)/n!$.
\end{definition}

\begin{definition}[Effective divisor]
  \label{def:divisor}
  \uses{def:disk}
  \lean{Weierstrass.EffectiveDivisor}
  \leanok
  An \emph{effective divisor} on $\D$ is a locally finite formal sum
  $D = \sum_{a \in S} m_a [a]$ with non-negative integer multiplicities
  over a discrete subset $S \subset \D$. Formally, $D$ is represented by
  its multiplicity function $\mathrm{mult} : \C \to \N$, vanishing outside
  $\D$, such that every compact $K \subseteq \D$ meets
  $\{z : \mathrm{mult}(z) \ne 0\}$ in only finitely many points (the
  standard characterization of a discrete, locally finite subset of the
  open set $\D$).
\end{definition}

\begin{definition}[Conjugation invariance]
  \label{def:conjInvariant}
  \uses{def:divisor}
  \lean{Weierstrass.EffectiveDivisor.ConjInvariant}
  \leanok
  An effective divisor $D$ is \emph{invariant under complex conjugation}
  if $m_{\bar a} = m_a$ for all $a$.
\end{definition}

\begin{definition}[Zero divisor of a function]
  \label{def:zeroDivisorOf}
  \uses{def:disk, def:divisor}
  \lean{Weierstrass.IsZeroDivisorOf}
  \leanok
  $D$ is \emph{the zero divisor} of $f$ if, for every $z \in \D$, the
  multiplicity $m_z$ equals the order of vanishing of $f$ at $z$.
\end{definition}

\chapter{Modified Elementary Factors}
\label{chap:factors}

\begin{definition}[Modified elementary factor]
  \label{def:E}
  \lean{Weierstrass.E}
  \leanok
  For $n \ge 0$ and $c \in \C$, the modified elementary factor of order
  $n$ with parameter $c$ is
  \[
    E_n(w;\,c) = (1-w)\exp\Bigl(\sum_{k=1}^{n}\frac{w^k}{k} +
      \frac{c\,w^{n+1}}{n+1}\Bigr).
  \]
  For $c = 1$ this is the classical Weierstrass elementary factor.
\end{definition}

\begin{definition}[The exponent $G_n$]
  \label{def:G}
  \lean{Weierstrass.G}
  \leanok
  \[
    G_n(w;\,c) = \frac{(c-1)\,w^{n+1}}{n+1} - \sum_{k=n+2}^{\infty}\frac{w^k}{k}.
  \]
\end{definition}

\begin{lemma}[$E_n = \exp(G_n)$ on $\D$]
  \label{lem:Eexp}
  \uses{def:disk, def:E, def:G}
  \lean{Weierstrass.E_eq_exp_G}
  \leanok
  For $w \in \D$, $E_n(w;\,c) = \exp(G_n(w;\,c))$.
\end{lemma}
\begin{proof}
  \uses{lem:Eexp}
  \leanok
  Since $-\log(1-w) = \sum_{k \ge 1} w^k/k$ for $w \in \D$, the exponent
  in the definition of $E_n$ equals $-\log(1-w)$ plus the affine
  correction term and minus the tail $\sum_{k \ge n+2} w^k/k$; regrouping
  and exponentiating cancels the $(1-w)^{-1}$ factor against $(1-w)$.
\end{proof}

\begin{lemma}[Structure of the modified factor]
  \label{lem:structure}
  \uses{def:disk, def:E, def:G, lem:Eexp}
  \lean{Weierstrass.E_zero, Weierstrass.taylorCoeff_E_eq_zero, Weierstrass.taylorCoeff_E_succ, Weierstrass.E_ne_zero, Weierstrass.E_zero_iff}
  \leanok
  For all $n \ge 0$, $c \in \C$, and $w \in \D$:
  \begin{enumerate}
    \item[(i)] $E_n(0;\,c) = 1$.
    \item[(ii)] $[w^m]\,E_n(\,\cdot\,;\,c) = 0$ for $1 \le m \le n$,
      independently of $c$.
    \item[(iii)] $[w^{n+1}]\,E_n(\,\cdot\,;\,c) = (c-1)/(n+1)$, affine in
      $c$ with slope $1/(n+1)$.
    \item[(iv)] $E_n(w;\,c) \ne 0$ for all $w \in \D$.
    \item[(v)] As an entire function of $w \in \C$, $E_n(\,\cdot\,;\,c)$
      has a simple zero at $w = 1$ and no other zeros.
  \end{enumerate}
\end{lemma}
\begin{proof}
  \uses{lem:structure}
  \leanok
  (i) is immediate from the definition. (ii)--(iii) follow from the power
  series expansion $G_n(w;\,c) = \tfrac{c-1}{n+1}w^{n+1} + O(w^{n+2})$
  together with Lemma~\ref{lem:Eexp}: composing with $\exp$ leaves
  degrees $\le n$ untouched and reads off the degree-$(n+1)$ coefficient.
  (iv) is immediate since $\exp$ is nowhere zero. (v) follows directly
  from the factorization $E_n(w;\,c) = (1-w)e^{H(w)}$ with $H$ entire, so
  the only zero is that of $(1-w)$; the order (simplicity) of this zero
  is established separately at Lemma~\ref{lem:zeroSet}, via
  \lean{Weierstrass.analyticOrderAt_E_div_self}.
\end{proof}

\chapter{Affine Coefficient Control}
\label{chap:affine}

\begin{lemma}[Affine coefficient control]
  \label{lem:affine}
  \uses{def:E, def:taylorCoeff}
  \lean{Weierstrass.taylorCoeff_mul_E_eq_of_le, Weierstrass.taylorCoeff_mul_E_succ}
  \leanok
  Let $a \in \C \setminus \{0\}$, $n \ge 0$, and $h$ a power series
  convergent near $0$ with $h_0 = 1$. For $F(z) = h(z)\cdot
  E_n(z/a;\,c)$:
  \begin{enumerate}
    \item[(i)] $[z^m]\,F = h_m$ exactly, for $0 \le m \le n$,
      independently of $c$.
    \item[(ii)] $[z^{n+1}]\,F = h_{n+1} + \dfrac{c-1}{(n+1)a^{n+1}}$,
      affine in $c$ with nonzero slope $1/((n+1)a^{n+1})$.
  \end{enumerate}
\end{lemma}
\begin{proof}
  \uses{lem:affine}
  \leanok
  Write $E_n(z/a;\,c) = \sum_j e_j z^j$; by Lemma~\ref{lem:structure},
  $e_0 = 1$, $e_1 = \dots = e_n = 0$, and
  $e_{n+1} = (c-1)/((n+1)a^{n+1})$. Then $[z^m]F = \sum_{j=0}^m
  h_{m-j}e_j$, and the vanishing of $e_1, \dots, e_n$ isolates the $j=0$
  term for $m \le n$, or the $j=0$ and $j=n+1$ terms for $m=n+1$.
\end{proof}

\begin{lemma}[Triangular-affine structure]
  \label{rem:triangular}
  \uses{lem:affine}
  \lean{Weierstrass.exists_c_taylorCoeff_mul_E_succ_eq}
  \leanok
  In the setting of Lemma~\ref{lem:affine}, the map $c \mapsto
  [z^{n+1}](h\cdot E_n(z/a;\,c))$ is a surjective affine map $\C \to \C$;
  in particular every target value is attained by some $c$.
\end{lemma}
\begin{proof}
  \uses{rem:triangular}
  \leanok
  Immediate from Lemma~\ref{lem:affine}(ii): the map is affine with
  nonzero slope $1/((n+1)a^{n+1})$, hence bijective.
\end{proof}

\chapter{Gaussian-Integer Realization}
\label{chap:propZi}

\begin{definition}[Gaussian-integer Taylor coefficients]
  \label{def:hasGaussianIntCoeffs}
  \uses{def:taylorCoeff}
  \lean{Weierstrass.HasGaussianIntCoeffs}
  \leanok
  $f$ \emph{has Taylor coefficients in $\Zi$} if $[z^n]f \in \Zi$ for
  every $n \ge 0$.
\end{definition}

\begin{theorem}[Gaussian-integer realization]
  \label{thm:propZi}
  \uses{def:disk, def:holomorphicOn, def:divisor, def:zeroDivisorOf,
    def:hasGaussianIntCoeffs}
  \lean{Weierstrass.exists_holomorphic_gaussianInt_coeffs_of_effectiveDivisor}
  \leanok
  Every effective divisor on $\D$ is the zero divisor of a holomorphic
  function on $\D$ with Taylor coefficients in $\Zi$.
\end{theorem}

The proof enumerates the divisor's support with multiplicity, then
inductively rounds the Taylor coefficients of the partial products to
the nearest Gaussian integer, using the triangular structure of
Lemma~\ref{lem:affine} so that rounding a new coefficient never disturbs
an earlier one. The origin is handled separately by an explicit monomial
factor $z^{m_0}$, so the construction proper enumerates only the support
away from $0$.

\begin{lemma}[Nearest Gaussian integer]
  \label{lem:nearestGaussianInt}
  \lean{Weierstrass.nearestGaussianInt, Weierstrass.norm_sub_nearestGaussianInt_le}
  \leanok
  Rounding the real and imaginary parts of $v \in \C$ independently to
  the nearest integer gives a Gaussian integer $T$ with
  $\abs{T - v} \le \sqrt2/2$.
\end{lemma}
\begin{proof}
  \uses{lem:nearestGaussianInt}
  \leanok
  Each coordinate's rounding error is at most $1/2$; combine via
  $\abs{\cdot} \le \sqrt2 \max(\abs{\cdot_{\mathrm{re}}},
  \abs{\cdot_{\mathrm{im}}})$.
\end{proof}

\begin{lemma}[Bound on the exponent $G_n$]
  \label{lem:normG}
  \uses{def:G}
  \lean{Weierstrass.norm_G_le}
  \leanok
  For $\norm{w} \le \rho < 1$,
  \[
    \norm{G_n(w;\,c)} \le \frac{\norm{c-1}}{n+1}\,\rho^{n+1}
      + \frac{\rho^{n+2}}{1-\rho}.
  \]
\end{lemma}
\begin{proof}
  \uses{lem:normG}
  \leanok
  Bound the affine term by the triangle inequality, and the geometric
  tail $\sum_{k \ge n+2} w^k/k$ termwise by $\norm{w}^k \le \rho^k$,
  summing the resulting geometric series.
\end{proof}

\begin{lemma}[Divisor enumeration]
  \label{lem:enum}
  \uses{def:disk, def:divisor}
  \lean{Weierstrass.exists_enum_of_effectiveDivisor}
  \leanok
  The support of an effective divisor $D$, away from the origin, can be
  enumerated as a sequence $a : \N \to \C$ realizing the multiplicity
  function by counting repeats, with $a_k \ne 0$ for all $k$ and
  $\norm{a_k} \to 1$ (any fixed $s<1$ bounds $\norm{a_k}$ away from $s$
  for all but finitely many $k$).
\end{lemma}
\begin{proof}
  \uses{lem:enum}
  \leanok
  $\D \setminus \{0\}$ is exhausted by compact annuli; on each, local
  finiteness of the divisor's support gives a finite list of points with
  multiplicity, and concatenating these lists over the annuli gives the
  required enumeration. (If the support is finite or empty, the
  enumeration is padded by a fixed point outside $\overline\D$, which
  contributes no zero in $\D$; this uniformly subsumes the paper's
  separate treatment of the finite-support case.)
\end{proof}

\begin{lemma}[Inductive coefficient forcing]
  \label{lem:coeffSeq}
  \uses{lem:affine, rem:triangular, lem:nearestGaussianInt, lem:enum}
  \lean{Weierstrass.exists_coeffSeq}
  \leanok
  Given $a : \N \to \C$ with $a_k \ne 0$, there is a sequence of
  correction constants $c : \N \to \C$ such that, writing $P_N :=
  \prod_{k<N} E_k(z/a_k;\,c_k)$:
  \begin{itemize}
    \item every Taylor coefficient of $P_N$ at degree $N$ lies in $\Zi$;
    \item the rounding error is controlled:
      $\norm{c_k - 1} \le (\sqrt2/2)(k+1)\norm{a_k}^{k+1}$.
  \end{itemize}
\end{lemma}
\begin{proof}
  \uses{lem:coeffSeq}
  \leanok
  Induction on $N$, using Remark~\ref{rem:triangular} to force
  $[z^{N+1}]P_{N+1}$ to the Gaussian integer nearest its pre-forced
  value (Lemma~\ref{lem:nearestGaussianInt}); Lemma~\ref{lem:affine}(i)
  guarantees that later factors never disturb an already-forced
  coefficient. Built in Lean by a structural recursion computing the
  partial product and the next correction constant simultaneously.
\end{proof}

\begin{lemma}[Weierstrass $M$-test for the product]
  \label{lem:Mtest}
  \uses{def:disk, lem:normG, lem:coeffSeq, lem:enum}
  \lean{Weierstrass.exists_Mtest_of_coeffSeq}
  \leanok
  Given the quantitative bound on $c$ from Lemma~\ref{lem:coeffSeq} and
  the escape property from Lemma~\ref{lem:enum}, the factors
  $E_k(z/a_k;\,c_k)$ satisfy a Weierstrass $M$-test on every compact
  $K \subseteq \D$: there is a summable bound $u_k$, uniform in $z \in K$.
\end{lemma}
\begin{proof}
  \uses{lem:Mtest}
  \leanok
  Fix $K \subseteq \D$ compact, contained in $\{\norm z \le r\}$ for some
  $r<1$, and $s$ with $r < s < 1$. Beyond finitely many $k$,
  $\norm{a_k} \ge s$, so $\norm{z/a_k} \le r/s < 1$ and
  Lemma~\ref{lem:normG} applies with $\rho := r/s$; the quantitative
  bound on $c_k$ makes the affine term collapse to
  $(\sqrt2/2)r^{k+1}$ (the $\norm{a_k}^{k+1}$ factors cancel exactly),
  and both resulting geometric series are summable. Converting the
  $G$-bound to a bound on $\norm{E_k - 1}$ uses
  $\norm{e^x - 1} \le 2\norm x$ for $\norm x \le 1$. The finitely many
  exceptional small $k$ are absorbed into the bound without affecting
  summability.
\end{proof}

\begin{lemma}[Holomorphy of the infinite product]
  \label{lem:holomorphicProduct}
  \uses{def:disk, def:holomorphicOn, lem:Mtest}
  \lean{Weierstrass.hasProdLocallyUniformlyOn_factors, Weierstrass.holomorphicOn_tprod_factors}
  \leanok
  Under an $M$-test bound as in Lemma~\ref{lem:Mtest} (stated for a
  general order function $n : \N \to \N$ and points/parameters
  $a, c : \N \to \C$), the partial products of
  $E_{n(k)}(z/a_k;\,c_k)$ converge locally uniformly on $\D$, and the
  limit is holomorphic on $\D$.
\end{lemma}
\begin{proof}
  \uses{lem:holomorphicProduct}
  \leanok
  Standard consequence of the Weierstrass $M$-test for products
  (\cite[Ch.~5, Thm.~1.1]{Conway}): locally uniform convergence of
  holomorphic partial products has a holomorphic limit.
\end{proof}

\begin{lemma}[Freezing of coefficients]
  \label{lem:taylorStabilize}
  \uses{lem:affine, lem:holomorphicProduct}
  \lean{Weierstrass.taylorCoeff_tprod_factors_eq_partial}
  \leanok
  For a monotone order function $n : \N \to \N$, once $n(k) > m$ the
  $m$-th Taylor coefficient of the infinite product
  $\prod'_j E_{n(j)}(z/a_j;\,c_j)$ agrees with that of the
  $(k+1)$-st partial product.
\end{lemma}
\begin{proof}
  \uses{lem:taylorStabilize}
  \leanok
  Lemma~\ref{lem:affine}(i) shows the finite partial products stabilize
  at degree $m$ once their length exceeds $k$; passing to the locally
  uniform limit, iterated derivatives (hence Taylor coefficients) pass
  to the limit along the same stabilized sequence.
\end{proof}

\begin{lemma}[Zero-divisor identification]
  \label{lem:zeroSet}
  \uses{lem:structure, lem:holomorphicProduct, lem:Mtest}
  \lean{Weierstrass.analyticOrderAt_tprod_factors_eq_ncard, Weierstrass.isZeroDivisorOf_tprod_factors, Weierstrass.analyticOrderAt_E_div_self}
  \leanok
  At every point $z \in \D$, the order of vanishing of the infinite
  product $\prod'_k E_{n(k)}(z/a_k;\,c_k)$ equals $\abs{\{k : a_k = z\}}$
  (counted with multiplicity, i.e.\ as a cardinality of a fiber).
\end{lemma}
\begin{proof}
  \uses{lem:zeroSet}
  \leanok
  By Lemma~\ref{lem:structure}(v), $E_n(\,\cdot/a;\,c)$ has a simple zero
  exactly at $a$; additivity of the order of vanishing under finite
  products gives the order of a partial product as the number of
  vanishing factors. Splitting the infinite product at an index beyond
  which no factor vanishes at $z$, the tail is holomorphic and
  non-vanishing near $z$ (by the same $M$-test criterion used for
  Lemma~\ref{lem:holomorphicProduct}), so it contributes order $0$.
\end{proof}

\begin{proof}[Proof of Theorem~\ref{thm:propZi}]
  \uses{thm:propZi, lem:enum, lem:coeffSeq, lem:Mtest,
    lem:holomorphicProduct, lem:taylorStabilize, lem:zeroSet}
  \leanok
  Factor out the origin: writing $D = m_0[0] + D'$, it suffices to
  realize $D'$ (supported away from $0$) by $g \in \OD$ with
  $\Zi$-coefficients and $g(0) \ne 0$, and set $f(z) = z^{m_0}g(z)$.
  Enumerate $D'$'s support via Lemma~\ref{lem:enum}, force the Taylor
  coefficients of the partial products into $\Zi$ via
  Lemma~\ref{lem:coeffSeq} (taking the order function $n = \mathrm{id}$,
  i.e.\ one factor per index, densely), and pass to the infinite
  product $g := \prod'_k E_k(z/a_k;\,c_k)$. Holomorphy is
  Lemma~\ref{lem:holomorphicProduct}; every Taylor coefficient of $g$
  lies in $\Zi$ by Lemma~\ref{lem:taylorStabilize} (each degree $m$
  stabilizes, by the $m+1$-st stage, to the value forced in
  Lemma~\ref{lem:coeffSeq}); and $g$'s zero divisor is $D'$ by
  Lemma~\ref{lem:zeroSet} together with the multiplicity guarantee of
  Lemma~\ref{lem:enum}.
\end{proof}

\chapter{The Main Theorem}
\label{chap:main}

\begin{definition}[Integer Taylor coefficients]
  \label{def:hasIntCoeffs}
  \uses{def:taylorCoeff}
  \lean{Weierstrass.HasIntCoeffs}
  \leanok
  $f$ \emph{has Taylor coefficients in $\Z$} if $[z^n]f \in \Z$ for
  every $n \ge 0$.
\end{definition}

\begin{theorem}[Main theorem]
  \label{thm:main}
  \uses{def:disk, def:holomorphicOn, def:divisor, def:zeroDivisorOf,
    def:conjInvariant, def:hasIntCoeffs}
  \lean{Weierstrass.exists_holomorphic_int_coeffs_iff_conjInvariant}
  \leanok
  An effective divisor $D$ on $\D$ is the zero divisor of a holomorphic
  function on $\D$ with Taylor coefficients in $\Z$ if and only if $D$
  is invariant under complex conjugation.
\end{theorem}

\section{Necessity}

\begin{lemma}[Double conjugation is holomorphic]
  \label{lem:conjComp}
  \lean{Weierstrass.hasDerivAt_conj_comp_conj, Weierstrass.analyticAt_conj_comp_conj}
  \leanok
  If $f$ has derivative $f'$ at $w$, then $z \mapsto \overline{f(\bar z)}$
  has derivative $\overline{f'}$ at $\bar w$. Consequently, if $f$ is
  analytic at $p$, then $z \mapsto \overline{f(\bar z)}$ is analytic at
  $\bar p$.
\end{lemma}
\begin{proof}
  \uses{lem:conjComp}
  \leanok
  A single conjugation (e.g.\ $f \circ \mathrm{conj}$) is generally
  anti-holomorphic, but doubly conjugating both the argument and the
  value is holomorphic again: this is checked at the level of the
  Fr\'echet derivative, using that conjugation is $\R$-linear and that
  conjugating a $\C$-linear scalar-multiplication map on both sides
  yields another $\C$-linear scalar-multiplication map, by the
  conjugated scalar.
\end{proof}

\begin{lemma}[Conjugation of Taylor coefficients and order]
  \label{lem:taylorCoeffConj}
  \uses{lem:conjComp, def:taylorCoeff}
  \lean{Weierstrass.taylorCoeff_conj_comp_conj, Weierstrass.analyticOrderAt_conj_comp_conj}
  \leanok
  If $f$ is analytic at $0$ (resp. at $p$), then
  $[z^n](\overline{f(\bar z)}) = \overline{[z^n]f}$ for all $n$, and the
  order of vanishing of $\overline{f(\bar z)}$ at $\bar p$ equals the
  order of vanishing of $f$ at $p$.
\end{lemma}
\begin{proof}
  \uses{lem:taylorCoeffConj}
  \leanok
  The Taylor-coefficient statement follows by iterating
  Lemma~\ref{lem:conjComp} order by order (using that eventual equality
  of two functions near a point transfers through iterated derivatives).
  The order statement follows by writing $f$ near $p$ as
  $(z-p)^k \cdot \varphi(z)$ with $\varphi$ analytic and non-vanishing at
  $p$, and doubly conjugating both $f$ and $\varphi$.
\end{proof}

\begin{lemma}[Necessity direction]
  \label{lem:necessity}
  \uses{def:disk, def:holomorphicOn, def:divisor, def:zeroDivisorOf,
    def:conjInvariant, def:hasIntCoeffs, lem:taylorCoeffConj}
  \lean{Weierstrass.conjInvariant_of_hasIntCoeffs}
  \leanok
  If $f$ is holomorphic on $\D$ with the zero divisor of $D$ and integer
  Taylor coefficients, then $D$ is invariant under complex conjugation.
\end{lemma}
\begin{proof}
  \uses{lem:necessity}
  \leanok
  Since every Taylor coefficient of $f$ is real (being an integer),
  Lemma~\ref{lem:taylorCoeffConj} shows $f$ and $z \mapsto
  \overline{f(\bar z)}$ have the same Taylor series at $0$; since both
  are analytic and $\D$ is connected, the identity theorem for analytic
  functions upgrades this to equality throughout $\D$. Applying the
  order-preservation half of Lemma~\ref{lem:taylorCoeffConj} at each
  $z \in \D$ then gives $m_{\bar z} = m_z$.
\end{proof}

\section{Sufficiency: the conjugate-pairing construction}

The sufficiency direction cannot simply invoke
Theorem~\ref{thm:propZi} as a black box: pairing a Gaussian-integer
realization $g$ with its reflection $\overline{g(\bar z)}$ produces real,
but not generally integer, coefficients (Remark on alternative
symmetrization in the paper). Instead, the construction of
Theorem~\ref{thm:propZi} is redone with the zeros of a conjugate-invariant
divisor grouped into \emph{slots}: a single slot for each real zero, and a
slot shared by each conjugate pair of nonreal zeros, indexed so that slot
$j$ uses factors of order $j$ at flat positions $2j$ and $2j+1$. Because
both factors of a slot share the same order, the correction constants
can be linked ($\bar c$ for the conjugate-pair companion) so that the
slot's combined contribution to the Taylor coefficients is always real,
enabling rounding to the nearest \emph{integer} rather than nearest
Gaussian integer.

\begin{lemma}[Paired enumeration of a conjugate-invariant divisor]
  \label{lem:pairedEnum}
  \uses{def:divisor, def:conjInvariant, lem:enum}
  \lean{Weierstrass.exists_pairedEnum_of_conjInvariant}
  \leanok
  Every effective divisor $D$ invariant under complex conjugation can be
  enumerated in slots of two: even positions $a_{2j}$ list one point per
  real zero or per conjugate-pair representative (with multiplicity),
  and odd positions carry the real/conjugate companion of the preceding
  even position (a fixed dummy point outside $\D$ when $a_{2j}$ is real,
  since no companion is needed).
\end{lemma}
\begin{proof}
  \uses{lem:pairedEnum}
  \leanok
  Apply Lemma~\ref{lem:enum} to the ``upper half'' divisor obtained by
  restricting $D$'s multiplicities to points with non-negative imaginary
  part (excluding $0$); since $D$ is conjugate-invariant, its full
  support is recovered from this upper half together with its complex
  conjugate.
\end{proof}

\begin{lemma}[Conjugate symmetry of a slot factor]
  \label{lem:slotConjSymmetric}
  \uses{def:E}
  \lean{Weierstrass.E_conj, Weierstrass.conjSymmetric_E_real, Weierstrass.conjSymmetric_E_pair}
  \leanok
  Conjugating $E_n(w;\,c)$ conjugates both parameters:
  $\overline{E_n(w;\,c)} = E_n(\bar w;\,\bar c)$ (since $E_n$'s defining
  formula involves only $+,-,*,/,\exp$ applied to rational numeric
  constants, all of which commute with conjugation). Consequently:
  \begin{itemize}
    \item if $a, c$ are real, $z \mapsto E_n(z/a;\,c)$ satisfies
      $\overline{h(\bar z)} = h(z)$ for all $z$;
    \item for any $a, c$, $z \mapsto E_n(z/a;\,c)\cdot E_n(z/\bar
      a;\,\bar c)$ satisfies $\overline{h(\bar z)} = h(z)$ for all $z$.
  \end{itemize}
\end{lemma}
\begin{proof}
  \uses{lem:slotConjSymmetric}
  \leanok
  Direct computation from $\overline{E_n(w;\,c)} = E_n(\bar w;\,\bar c)$.
  A function satisfying $\overline{h(\bar z)} = h(z)$ everywhere and
  analytic at $0$ has real Taylor coefficients there, by
  Lemma~\ref{lem:taylorCoeffConj}
  (\lean{Weierstrass.taylorCoeff_real_of_conjSymmetric}); products of
  such functions again satisfy the identity.
\end{proof}

\begin{lemma}[Paired coefficient forcing]
  \label{lem:pairCoeffSeq}
  \uses{lem:affine, lem:pairedEnum, lem:slotConjSymmetric}
  \lean{Weierstrass.chooseCReal, Weierstrass.chooseCPair, Weierstrass.exists_pairCoeffSeq}
  \leanok
  Given the paired enumeration $a$ of Lemma~\ref{lem:pairedEnum}, there
  is a correction-constant sequence $c : \N \to \C$ such that:
  \begin{itemize}
    \item at each slot $j$, the degree-$(j+1)$ Taylor coefficient of the
      length-$2(j+1)$ partial product lies in $\Z$;
    \item $\norm{c_k - 1} \le (\sqrt2/2)(\lfloor k/2\rfloor+1)
      \norm{a_k}^{\lfloor k/2\rfloor + 1}$ for all $k$.
  \end{itemize}
  For a real-zero slot ($a_{2j}$ real), $c_{2j}$ is a real correction
  constant forcing the coefficient to the nearest integer to its
  (already-real, by Lemma~\ref{lem:slotConjSymmetric}) pre-existing
  value, with rounding error $\le 1/2$; $c_{2j+1} = 1$ is trivial. For a
  conjugate-pair slot, $c_{2j+1} = \overline{c_{2j}}$, and $c_{2j}$ is
  chosen as the minimal-norm complex number forcing the combined,
  real-affine shift $2\,\mathrm{Re}\bigl((c-1)/((j+1)a^{j+1})\bigr)$ to
  the nearest integer, giving rounding error $\le (j+1)/4 \cdot
  \norm{a_{2j}}^{j+1}$.
\end{lemma}
\begin{proof}
  \uses{lem:pairCoeffSeq}
  \leanok
  For the real slot: identical to the inductive step of
  Lemma~\ref{lem:coeffSeq}, using Lemma~\ref{lem:affine} with a real
  target and real $a$, giving a real correction constant automatically
  (not as an extra constraint). For the pair slot: applying
  Lemma~\ref{lem:affine}(ii) twice (once for each factor of the pair)
  gives a total shift $\tfrac{1}{j+1}\bigl((c-1)/a^{j+1} +
  (\bar c - 1)/\bar a^{j+1}\bigr) = \tfrac{2}{j+1}\mathrm{Re}((c-1)/a^{j+1})$;
  writing $c - 1 = K \cdot a^{j+1}$ for real $K$ solves
  $\mathrm{Re}((c-1)\overline{a^{j+1}}) = K\norm{a}^{2(j+1)}$ with minimal
  $\norm{c-1} = \abs K \norm a^{j+1}$, and choosing the real integer
  target for $K$ (up to the factor $(j+1)/2$) gives the stated bound via
  the standard nearest-integer estimate $\abs{\mathrm{round}(x)-x}\le
  1/2$. A structural recursion (paralleling Lemma~\ref{lem:coeffSeq}'s)
  builds the partial products and correction constants simultaneously,
  slot by slot, maintaining the conjugate-symmetry invariant of
  Lemma~\ref{lem:slotConjSymmetric} throughout.
\end{proof}

\begin{lemma}[Paired Weierstrass $M$-test]
  \label{lem:pairMtest}
  \uses{def:disk, lem:normG, lem:pairCoeffSeq, lem:pairedEnum}
  \lean{Weierstrass.exists_Mtest_of_pairCoeffSeq}
  \leanok
  Given the bound on $c$ from Lemma~\ref{lem:pairCoeffSeq}, the factors
  $E_{\lfloor k/2\rfloor}(z/a_k;\,c_k)$ satisfy a Weierstrass $M$-test on
  every compact $K \subseteq \D$.
\end{lemma}
\begin{proof}
  \uses{lem:pairMtest}
  \leanok
  As in Lemma~\ref{lem:Mtest}, but with order function
  $n(k) = \lfloor k/2\rfloor$ in place of $n(k)=k$: since
  $\lfloor k/2 \rfloor \to \infty$ at half the rate, the geometric
  majorant $r^{\lfloor k/2\rfloor+1}$ is compared against $(\sqrt r)^k$
  (using $2\lfloor k/2\rfloor + 2 \ge k$ for $\sqrt r < 1$), which is
  still summable.
\end{proof}

\begin{proof}[Proof of Theorem~\ref{thm:main}, sufficiency direction]
  \label{lem:sufficiency}
  \uses{thm:main, lem:pairedEnum, lem:pairCoeffSeq, lem:pairMtest,
    lem:holomorphicProduct, lem:taylorStabilize, lem:zeroSet}
  \lean{Weierstrass.exists_holomorphic_int_coeffs_of_conjInvariant}
  \leanok
  Enumerate $D$ via Lemma~\ref{lem:pairedEnum} and force integer
  coefficients slot by slot via Lemma~\ref{lem:pairCoeffSeq}, with order
  function $n(k) = \lfloor k/2\rfloor$. The general holomorphy, Taylor
  freezing, and zero-set lemmas of Section~\ref{chap:propZi}
  (Lemmas~\ref{lem:holomorphicProduct}, \ref{lem:taylorStabilize},
  \ref{lem:zeroSet}) are stated for a general (monotone, not
  necessarily injective) order function, so they apply verbatim with
  $n(k) = \lfloor k/2 \rfloor$ and the $M$-test of
  Lemma~\ref{lem:pairMtest} in place of Lemma~\ref{lem:Mtest}. As in the
  proof of Theorem~\ref{thm:propZi}, the origin is handled by an
  explicit monomial factor $z^{m_0}$.
\end{proof}
